Siguin els plans $\pi_1$ i $\pi_2$, determinats respectivament per les equacions $\pi_1\colon x+y=3$ i $\pi_2\colon x-z=-2$.

\begin{apartats}

\apartat{0,75}
Trobeu l'equació general ($Ax+By+Cz+D=0$) del pla $\pi_3$, que és perpendicular a $\pi_1$ i $\pi_2$, i que passa pel punt
$P=(4,1,2)$.

\begin{solucio}
Com que el pla que busquem és perpendicular a $\pi_1$ i a $\pi_2$, els seus vectors directors seran els vectors normals
de $\pi_1$ i $\pi_2$: $\vec v_1=(1,1,0)$ i $\vec v_2=(1,0,-1)$. Trobem l'equació general del pla que té vectors directors
$\vec v_1$ i $\vec v_2$ i que passa pel punt $P=(4,1,2)$, amb $X(x,y,z)$ un punt genèric del pla:
\[
\begin{vmatrix}x-4&1&1\\y-1&1&0\\z-2&0&-1\end{vmatrix}=0\;\rightarrow\;\boxed{x-y+z-5=0}.
\]

\textit{Pauta oficial:} 0,25 pels vectors directors, 0,25 per indicar com trobar el pla i 0,25 pel càlcul de l'equació.
\end{solucio}

\apartat{0,75}
Sigui $r$ la recta d'intersecció de $\pi_1$ i $\pi_2$. Calculeu l'equació vectorial de la recta $r$.

\begin{solucio}
Busquem la solució del sistema $\left\{\begin{aligned}x+y&=3\\x-z&=-2\end{aligned}\right.$. Aïllant, obtenim $y=3-x$ i
$z=2+x$, i, si $x=\lambda$, tenim $x=\lambda$, $y=3-\lambda$ i $z=2+\lambda$. La recta buscada, en forma vectorial, és
\[
\boxed{r\colon(x,y,z)=(0,3,2)+\lambda\cdot(1,-1,1),\ \lambda\in\mathbb{R}}.
\]

\textit{Pauta oficial:} 0,25 per formular la resolució del sistema i 0,5 per trobar l'equació.
\end{solucio}

\apartat{1}
Calculeu el punt $Q$ de la recta $r$ que és més a prop del punt $P$.

\begin{solucio}
Trobem l'equació del pla perpendicular a la recta $r$ que passa pel punt $P$. Aquest és el pla que ens han demanat a
l'apartat a), $x-y+z-5=0$. (En cas que no ens adonem que ja el tenim, aquest pla tindrà com a vector normal el vector
director de la recta, i serà de la forma $x-y+z+D=0$; imposant que $P$ hi sigui, $4-1+2+D=0\rightarrow D=-5$.) Ara hem
de calcular el punt de tall de la recta $r$ i el pla: substituint $x=\lambda$, $y=3-\lambda$, $z=2+\lambda$ a
$x-y+z=5$, obtenim $\lambda-(3-\lambda)+(2+\lambda)=5\rightarrow\lambda=2$, i, substituint $\lambda=2$ a les equacions
paramètriques de la recta, obtenim el punt $\boxed{Q=(2,1,4)}$.\\
De forma alternativa: el punt $Q$ és la projecció ortogonal del punt $P$ sobre la recta $r$, ja que és el punt que està a
menor distància de $P$. Sigui $Q=(\lambda,\,3-\lambda,\,2+\lambda)$ un punt genèric de la recta $r$, i $\vec v=(1,-1,1)$
el vector director de $r$. Per la condició d'ortogonalitat, tenim que $\overrightarrow{PQ}\cdot\vec v=0$:
\[
(-4+\lambda,\,2-\lambda,\,\lambda)\cdot(1,-1,1)=0\;\rightarrow\;-4+\lambda-2+\lambda+\lambda=0\;\rightarrow\;\lambda=2,
\]
i, substituint $\lambda=2$ a la recta $r$, obtenim que $Q=(2,1,4)$.

\textit{Pauta oficial:} 0,25 per indicar que el punt buscat és la projecció ortogonal, 0,25 pel punt genèric (o per
plantejar el sistema entre el pla i la recta), 0,25 pel producte escalar (o per resoldre el sistema) i 0,25 per trobar
el punt.
\end{solucio}

\end{apartats}
